\documentclass[11pt,letterpaper]{article}
\usepackage{styles/michaillat-paper}
\usepackage{styles/working-paper}
\usepackage{styles/yax-common}
\geometry{margin=1in}
\hypersetup{pdftitle={Response to Two Referee Reports},pdfauthor={Lily Luo}}
\setlength{\parindent}{0pt}
\setlength{\parskip}{0.55em}
\begin{document}
\begin{center}
{\Large\bfseries Response to Two Referee Reports}\par
\vspace{0.35em}
{\large ``Occupational AI Exposure and Young-Worker Employment:\\
Support and Comparisons in the CPS''}\par
\vspace{0.35em}
Lily Luo \quad September 2026
\end{center}

\section*{Response to the editor}

We thank both referees for identifying the same central problem from different directions: the earlier manuscript reported a sizeable sensitivity to broad-family conditioning, but it did not yet make occupational support and the precise descriptive target the organizing contribution.  We have substantially refocused the paper.  Its positive result is now that an apparently broad national exposure comparison contains little direct within-family tail information.  On the primary 468-occupation support, only four broad families contain both Q1 and Q5 occupations; the resulting 29 direct-tail occupations represent 5.03 percent of preperiod employment stock, and their information is equivalent to 6.2 occupations.  The paper treats the pooled-to-conditioned coefficient movement as a consequence of that support geometry, not as a causal ``share explained'' or evidence against an AI effect.

The revision adds the analyses required to make that claim auditable.  It places the family-by-quintile support matrix and information diagnostics in the main paper; reports 89 supported family-specific edges and a broader 490-occupation support check; gives exact young/older and within-family/family-weight stock accounting; crosses computer-use and family conditioning on common support; separates the static coefficient from reference-dependent and reference-invariant dynamic functionals; completes a target-specific finite-sample audit; adds age-band, fixed-age-composition, enrollment, cohort-composition, flow-selection, ACS, public-benchmark, and later-adoption checks; and resolves the earlier numerical-existence question for all eleven declared models.  The manuscript also distinguishes results that are informative under one sampling target from those that are robust across occupation, family, household, and ACS-replicate perturbations.

Several requests required modification rather than nominal compliance.  We do not promote December 2024 to the primary endpoint after observing the outcomes; it is instead a prominent pre-2025-production-break benchmark beside the unchanged full-window target.  We do not manufacture an Oster-style scalar bound for a nonlinear grouped-binomial coefficient when the required selection relation is not identified.  Exact BCC occupation memberships and tie handling are not public, so our BCC exercises remain labeled public-data analogues rather than replications.  Full-window weekly earnings remain unavailable because the authorized extracts omit \texttt{EARNWEEK2}; the short legacy series is retained only as ancillary evidence.  These non-adoptions and input constraints are stated below rather than described as completed work.

\section*{Response to Referee A: major comments}

\subsection*{3.1. Family conditioning, support, and the order of the argument}

\emph{Response: accepted.}  The abstract, introduction, and first results section now lead with support and information rather than the near-zero conditioned coefficient.  The main paper displays the family-by-quintile matrix and reports that only four of 22 families span Q1 and Q5, that the 29 direct-tail occupations account for 5.03 percent of preperiod stock, that family-month conditioning retains 29.7 percent of pooled Q5 target information, and that the direct-tail effective count is 6.2 occupations with a 77.8-percent top-five information share.  The direct-tail interval, $[-0.169,0.468]$, is presented before interpreting the conditioned coefficient.

The introduction now states explicitly that family paths cannot distinguish an AI effect operating at the broad-family level from unrelated family-level evolution.  The movement from $-0.1321$ to $-0.0217$ is therefore a comparison of conditional estimands and a corollary of the support geometry.  It is not a causal decomposition, and the residual coefficient is not evidence that AI has no effect.

\subsection*{3.2. Pandemic quarters and preperiod diagnostics}

\emph{Response: accepted, with the nulls kept distinct.}  The appendix now reports fixed tests for 2017Q1--2019Q4, 2021Q1--2022Q3, and the full preperiod excluding 2020Q2--2020Q4, under both the short-reference null and a reference-invariant within-window equality null.  Over the full preperiod, the reference-invariant test rejects for the pooled model ($p=.0121$) but not for the family-conditioned model ($p=.0641$).  After excluding 2020Q2--2020Q4, neither within-window equality test rejects at five percent ($p=.0676$ pooled and $p=.1048$ conditioned).  Relative to the short 2022Q4 reference, the corresponding pandemic-exclusion values are $p=.0772$ and $p=.0079$.  The difference is reported as a property of the null and the unusually short reference period, not as a favorable re-selection of the window.

Leave-one-quarter diagnostics identify 2019Q2 as producing the largest reduction in the pooled reference-relative Wald statistic and 2018Q4 in the conditioned statistic.  Thus the revised paper does not say that rejection is confined to pandemic quarters.  It says that the full pooled within-window rejection no longer clears five percent when those quarters are excluded, while the reference-relative family rejection persists.  HonestDiD is now run on the full event vector and covariance for the explicitly named post functional $P$, with both the full calibration and the consecutive 2021Q1--2022Q3 calibration.  Its conventional intervals already include zero, and the first positive smoothness allowance also includes zero; the exercise is not presented as a robustness statement for the static coefficient.

\subsection*{3.3. Young numerator and older denominator}

\emph{Response: addressed through exact stock accounting rather than unmatched regressions.}  Under the saturated occupation--month nuisance structure, separate single-age regressions do not reproduce the ratio target: the corresponding single-age target is absorbed, while an additive occupation-plus-month regression would define a different model.  We therefore report an exact, model-free equal-month accounting on the same Q1 and Q5 stocks.  The Q5--Q1 young-stock change is $-0.1098$ log point and the older-stock change is $0.0483$, so their young-relative difference is $-0.1581$.  The numerator supplies the larger movement and the denominator reinforces it.  A second exact Shapley accounting assigns $-0.1293$ to within-family ratio changes and $-0.0288$ to changing older-family weights, with a zero boundary term after auditing all absent cells.

These identities are reported in the main paper and appendix, but are not forced to equal the nonlinear grouped-binomial coefficient $-0.1321$.  We label them descriptive stock accounting, not separate causal regressions or percentage shares.

\subsection*{3.4. Finite-sample stress results}

\emph{Response: accepted.}  We withdrew the earlier 26.7- and 11.3-percent figures rather than defending them.  The replacement audit has eleven declared designs, 799--1,599 retained outer replications per design, 9,999 common multiplier draws per joint refit, independently computed pseudo-targets, and separate occupation, family, and cross-fit-oracle procedures.  In the empirical observed-effect design, occupation-wild coverage is 0.865, 0.921, and 0.889 for the pooled, family-conditioned, and paired targets; family-Rademacher coverage is 0.990, 0.895, and 0.992.  In the clean zero-target design, rejection of the family-conditioned target is 6.1 percent under occupation weights, 8.9 percent under family weights, and 5.6 percent for the cross-fit oracle.

The simulation also isolates why a structural-null calibration can remain far from a zero pooled projection: removing the structural exposure coefficient while retaining the fitted family-month surface leaves a pooled pseudo-target of $-0.11035$.  The revision therefore does not offer one omnibus ``size-corrected'' interval.  It reports target-specific performance and states the stronger limitation: no feasible cluster procedure is validated for every target and calibrated design, and rejection of the pooled projection alone is not a test of a structural exposure effect.

\subsection*{3.5. Computer use, collinearity, and sensitivity}

\emph{Response: accepted in substance; the proposed scalar bound was not imposed.}  On a fixed 455-occupation support, the paper now reports the four matched models: baseline, computer use, family-month paths, and both together, all with common draws.  Computer use moves the exposure coefficient from $-0.0957$ to $-0.1966$; family paths alone give $0.0354$; both give $-0.0467$.  The computer-use coefficient itself is $0.0460$ per preperiod-weighted standard deviation in the pooled computer model and $0.0487$ in the combined model, with both intervals reported.  The beta--computer-use correlation is 0.797 on this support, retained Q5 information falls from 69.1 to 29.0 percent, and the VIF-like ratio is 3.45.

The text treats the positive computer coefficient and the more negative exposure coefficient as joint-projection facts under collinearity, not as evidence that computerization raises employment or that the residual isolates AI.  We added an explicit applicability ruling for Oster, Cinelli--Hazlett, and Diegert--Masten--Poirier style analyses: none maps automatically to this nonlinear grouped-binomial target with generated, highly collinear treatments and changing nuisance spaces.  Rather than choose an arbitrary proportional-selection parameter after outcomes were known, we report the matched projection matrix and information loss and leave a scalar confounder bound unreported.

\subsection*{3.6. Enrollment and cohort composition}

\emph{Response: accepted.}  Restricting the young numerator to nonenrolled workers moves the pooled coefficient from $-0.1321$ to $-0.1239$, with paired movement $0.0082$ and interval $[-0.0125,0.0289]$.  On the common ages-22--54 question universe, restricting both groups to nonenrolled workers produces a similarly small paired movement.  The appendix reports annual enrollment and BA-or-higher shares by quintile, exact-age composition, approximate birth-year composition, and a national ages-22--25 population profile that assigns no exposure to nonworkers.  National enrollment changes from 0.2336 to 0.2211, BA-or-higher from 0.2941 to 0.3121, and average monthly weighted population from 16.94 to 17.60 million between the defined pre- and postperiods.  These results document cohort composition as a live alternative; they do not turn enrollment into a causal control.

\subsection*{3.7. Endpoints and survey-production changes}

\emph{Response: accepted on magnitudes; the proposed change of primary endpoint was not adopted.}  The endpoint discussion now compares movements rather than sorting nearly identical estimates by whether one interval endpoint crosses zero.  Relative to July 2026, the pooled movements through December 2024, September 2025, and December 2025 are 0.0211, 0.0225, and 0.0212 log point.  Their magnitudes are essentially the same even though their interval endpoints differ.

We retain the full available window as the canonical target because December 2024 would be selected after observing that it preserves a negative detected coefficient.  December 2024 is reported prominently as the pre-2025-production-break benchmark: its pooled coefficient is $-0.1110$ with interval $[-0.2008,-0.0213]$, and its movement from the full-window estimate has interval $[-0.0073,0.0495]$.  The paper separately reports official-weight and respondent-equivalent era comparisons, exclusions around the October 2025 gap, and the January 2025 and January 2026 control changes.  It does not claim that aggregate control factors can reconstruct unavailable age-by-occupation counterfactual weights.

\subsection*{3.8. Crosswalk allocation uncertainty}

\emph{Response: accepted as a finite sensitivity, not as a claimed bound.}  The current-contract analysis reallocates the 14.86 percent of supported early stock that passes through 31 one-to-many source occupations.  A seven-point jointly feasible odds-tilt grid preserves source--age--month mass and official structural zeros, and recomputes exposure groups after each allocation.  Pooled coefficients range from $-0.1446$ to $-0.1163$, with all six paired intervals relative to no tilt containing zero.  The family-month coefficient is less stable: at $K=0.25$ it moves from $-0.0217$ to $0.0439$, a paired change of $0.0656$ with interval $[0.0031,0.1280]$.  A wider 25-point feasible grid puts the pooled coefficient between $-0.1599$ and $-0.0702$.

Because the public files do not reveal true age-specific route probabilities, these are declared feasible-grid sensitivities rather than sharp bounds or an identified set.  The superseded historical-contract tilt is no longer substituted for the rebuilt-contract question.

\subsection*{3.9. Link attrition and flow selection}

\emph{Response: accepted within observable limits.}  The revision reports link retention by age, quintile, period, and horizon, plus observable differences between retained and eligible-unlinked origins.  It adds origin-observable poststratification and missing-outcome accounting intervals.  For example, the descriptive twelve-month exit contrast moves from 0.0166 under official link weights to 0.0210 after poststratification, but the maximum untrimmed factor is 102.7 and the procedure requires conditional missing at random.  Without that assumption, the jointly feasible linear accounting interval is $[-1.570,1.575]$.

These results show why clustering cannot cure selection and why the direction of unobserved attrition bias is not identified.  We do not assert attenuation toward zero or use the poststratified result as a correction.  Flow estimates remain low-resolution descriptions on their own linked risk sets.

\subsection*{3.10. Comparison with BCC public evidence}

\emph{Response: modified because exact membership is unavailable.}  We added a public benchmark section before the paper's main conditional models.  It distinguishes BCC's ADP stock-growth shortfall from CPS regression coefficients, and distinguishes equal-occupation public-dashboard grouping from employment-stock weighting.  Exact dashboard memberships and tie handling are not released, so our top-two/bottom-three rule is explicitly an approximation, not a replication.

In the CPS bridge, the approximation gives $-0.0720$ pooled and $-0.0168$ with family-month conditioning; the paired movement is 0.0553 with interval $[-0.0013,0.1118]$.  The less saturated family-post movement is narrowly detected, while the stronger family-month movement is not.  A separate all-employed public benchmark and an annual ACS reconstruction compare the published horizons and populations.  The ACS BCC analogues differ from the published values, but without exhaustive public membership the gap cannot be allocated among sample, grouping, and estimator.  The revised paper therefore promotes the bridge while preserving this identification limit.

\subsection*{3.11. ACS and adoption data}

\emph{Response: accepted through two scoped extensions.}  The annual ACS extension uses the 2017--2019 and 2021--2024 one-year PUMS files and all 80 replicate weights per endpoint year; no nonexistent standard 2020 one-year file is imputed.  It materially narrows person-sampling uncertainty, yet the occupation- and family-shock intervals for every family-year variant include zero and the direct tail support does not expand.  This distinguishes person-sampling resolution from occupational comparison support.

For realized use, we merge the fixed exposure score to the public Real-Time Population Survey detailed-occupation workbook.  The exact-code match retains 121 nonsuppressed occupations, with an unweighted exposure--adoption correlation of 0.504 and a Q5--Q1 mean-adoption difference of 0.267.  The workbook pools four 2025--26 waves and lacks detailed-cell standard errors and adoption timing, so we report no causal interaction or p-values.  We did not force the industry--state BTOS measure onto occupation-level CPS outcomes through an ecological bridge.  The extension establishes descriptive alignment between exposure and later use, not treatment assignment.

\section*{Response to Referee A: minor comments}

\begin{enumerate}[leftmargin=1.8em]
\item The correct count is 6,188,956 route-expanded descendants.  The erroneous 36,188,956 count is removed; 5,322,047 eligible source records and the distinction between fractional descendants and persons are stated beside it.
\item ``Chapter'' has been replaced by ``paper.''
\item The industry comparison is labeled as an industry-cell baseline with a different objective and is separated from occupation-level controls.
\item The family-path figure now uses a symmetric-log vertical scale, keeps literal zero-stock months visible, and states the linear and logarithmic regions in its note.
\item Every computer-use correlation and estimate names its support.  The matched four-model main comparison uses 455 occupations; the broader characteristic table separately identifies its 347-occupation common support.
\item Principal MDE quantities are consolidated in one appendix precision table with their scales and target types.  Main-text repetitions are limited to cases where resolution is itself the result.
\item The introduction now states first what the paper establishes---the support geometry and the comparisons it localizes---and then states the causal and equivalence limits.  Repeated implementation caveats have been moved to the appendix and archive.
\item The literature discussion now engages Bick, Blandin, and Deming; Hampole, Papanikolaou, Schmidt, and Seegmiller; Deming and Noray; and Autor and Thompson, and explains how adoption, firm investment, technological obsolescence, and expertise bear on the descriptive target.
\item The abstract defines the log conditional-mean young-to-older employment-stock ratio before reporting coefficient values and distinguishes it from an employment probability.
\item The conclusion now signposts Online Appendix I.3's consolidated data, design, and interpretation limits.
\end{enumerate}

\section*{Response to Referee B: major comments}

\subsection*{1. Benchmark and contribution}

\emph{Response: accepted.}  The revised contribution is not that a literal merge error proves a published finding wrong.  The introduction explicitly declines that attribution without another study's implementation evidence.  Instead, the contribution is that a national occupational AI-exposure comparison can contain much less direct within-family information than its nominal occupation count suggests.  The paper defines a public benchmark and separately records differences in population, outcome, grouping, weighting, endpoint, nuisance structure, and available employer controls.  The title and interpretation no longer claim to adjudicate ``AI versus composition.''  They describe support and comparisons and preserve the observational equivalence between family-level AI effects and family-level confounding.

\subsection*{2. From coefficient attenuation to explicit accounting}

\emph{Response: accepted.}  The revision distinguishes three objects: employment reallocations across families, within-family young-to-older ratio changes, and movements in a fitted coefficient after family paths are added.  Exact log-stock accounting gives a young Q5--Q1 change of $-0.1098$, an older change of $0.0483$, and a relative change of $-0.1581$.  A symmetric Shapley decomposition of the same family stocks assigns $-0.1293$ to within-family ratio changes and $-0.0288$ to changes in older-family weights.  It audits 36 structurally absent family-tail-period cells and obtains a zero boundary term.  This accounting is deliberately kept distinct from the nonlinear regression coefficient.

We also added the requested matched-support four-model comparison.  On the same 455 occupations and common draws, the baseline, computer-use, family-month, and combined coefficients are $-0.0957$, $-0.1966$, $0.0354$, and $-0.0467$.  Thus family and computer-use conditioning are different projections whose effects interact; no one residual is labeled an AI component.

\subsection*{3. Direct overlap and common-coefficient restrictions}

\emph{Response: accepted.}  The family support matrix is now in the main paper, alongside direct-tail coverage, information retention, effective occupation count, and concentration.  The text explains how intermediate quintiles and a common profile connect families that do not span Q1 and Q5, and notes that computer and mathematical occupations contribute only a small share of the conditioned target's information.

The appendix now retains all 89 supported family-specific pairwise edge functionals.  Their joint covariance has rank 50, so we do not report an invalid 89-degree-of-freedom Wald test or select a discovery list after inspecting pointwise results; the associated rank-50 coefficient system is reported with rank-aware inference.  Relaxing Webb availability expands support from 468 to 490 occupations and direct tails from 29 to 62 occupations and from 5.03 to 12.61 percent of stock.  SOC43 becomes a fifth spanning family, but overlap remains confined to five of 23 families.  The pooled coefficient changes little.  This shows that Webb availability contributes to, but does not create, the support limitation.

\subsection*{4. Static and dynamic reconciliation}

\emph{Response: accepted.}  We now report three targets with fixed weights and covariance-preserving comparisons.  The static grouped-binomial target $S$ is $-0.1321$ pooled and $-0.0217$ conditioned.  The post functional relative to the short 2022Q4 reference, $P$, is $-0.1199$ and $-0.2074$.  The reference-invariant post-minus-pre functional $D$ is $-0.1314$ and $-0.0217$.  The paired $D-S$ differences are 0.0007 ($p=.454$) and less than 0.0001 ($p=.995$), so the design does not detect differences between the static and reference-invariant dynamic targets.  By contrast, family conditioning changes $P-S$ by $-0.1980$ ($p=.021$).  This quantifies how the short reference period, rather than a contradiction in the post-minus-pre comparison, produces the apparent reversal.  HonestDiD is explicitly limited to $P$ and is not used to make claims about $S$ or $D$.

\subsection*{5. Inference, numerical existence, and stable-taxonomy checks}

\emph{Response: accepted.}  The finite-sample audit now evaluates occupation-Rademacher, occupation-Webb, family-Rademacher, family-Webb, and cross-fit full-refit oracle procedures for the pooled, family-conditioned, and paired targets under empirically motivated and adverse designs.  It reports target-specific coverage rather than treating one stress number as the CPS sampling law.  Household resampling remains explicitly a companion for repeated sampled units and released weights, not independent validation of the saturated family-month parameter.  Elapsed-calendar HAC is also kept as a distinct stochastic target.

The appendix now states existence results directly.  All eleven declared grouped-binomial models pass rank, recession-direction, independent trust-path versus zero-start Newton/IRLS agreement, fitted-mean, gradient, Hessian, and two-sided profile checks at unchanged thresholds.  Post-2020 fits profile one boundary nuisance group spanning 77 positive-total rows; occupation-seasonality fits profile 110 boundary groups spanning 770 rows.  The groups are handled on the extended-likelihood face rather than deleted, no direction can move the reported treatment target to infinity, and every focal target is finite.  L-BFGS-B remains a disclosed failed-score diagnostic, not the corroborating solver.

The coding-stable post-2020 pooled and family-month estimates are $-0.1181$ and $-0.0304$; their paired movements from the full-window fits have intervals $[-0.0288,0.0569]$ and $[-0.1053,0.0878]$.  Occupation-specific seasonality fits also pass and move the targets by at most 0.0011.  Earlier failed attempts remain in the registry but are no longer the only evidence on these questions.

\subsection*{6. Age comparison and links to flows}

\emph{Response: accepted, with remaining data limits explicit.}  The paper now compares ages 22--25 with 26--30, 31--40, 41--50, 51--65, and the original 26--65 group on fixed occupational support.  Pooled coefficients remain negative, and all family-month coefficient intervals contain zero.  Holding the older group's exact-age composition fixed moves the pooled target by only 0.0008 with interval $[-0.0011,0.0027]$.  These are descriptive denominator checks; no older band is called untreated.

The flow appendix reports weighted baseline transitions, retention by age, quintile, period, and horizon, observable differences between linked and eligible-unlinked origins, poststratification, and missing-outcome accounting intervals.  Entry destinations are also placed on one common nonemployment denominator, including continued nonemployment and destinations outside the support, while the conditional Q1--Q5 allocation remains separately labeled.  All primary flow intervals include zero, so the paper shortens their interpretation and does not infer an absent mechanism.  A complete stock--flow calibration remains infeasible because aging, all job-to-job moves, unsupported destinations, denominator flows, and employer matches are not jointly observed on one risk set.

The short weekly-earnings series through March 2023 is retained only as an ancillary, employment-conditioned result.  A full-window earnings analysis is not claimed: \texttt{EARNWEEK2} is absent from the authorized extracts.

\section*{Response to Referee B: specific presentation and reproducibility comments}

\begin{enumerate}[leftmargin=1.8em]
\item \textbf{Dependent variable.}  Main Table 2 now says that the grouped-binomial criterion is estimated on the two age-specific weighted stocks and that coefficients parameterize log ratios of conditional means.  No observed log ratio is constructed, and one-sided zero cells remain.
\item \textbf{Exposure-group reconstruction.}  Characteristic tables state whether groups are rebuilt on a support or fixed from the primary universe.  The central computer/family comparison uses fixed assignments, one 455-occupation support, one calendar, and common draws.  Maximal-support and 347-occupation common-support rows are separated, with their different baselines explained in the notes.
\item \textbf{Qualifications and implementation history.}  The main paper is reorganized around support, accounting, and a smaller set of economic targets.  Repeated caveats, installation history, superseded numerical results, hashes, and failure receipts are moved to the online appendix or archive.  Nondetection-versus-equivalence and association-versus-causation rules are stated once in the design section and applied consistently.
\item \textbf{Reproducibility and submission fields.}  Central exhibits are tied to named treatment, population, estimator, and inference objects, with a machine-readable results ledger and aggregate-only exhibit checks.  This is an auditable package, not a claim that the referees independently replicated it.  Affiliation, email, acknowledgments, funding, conflicts, and journal disclosures remain visibly marked \texttt{[AUTHOR TO COMPLETE]}; none has been invented.
\end{enumerate}

\section*{Claims retired and unresolved inputs}

The paper no longer claims a causal AI effect, a causal composition share, economic equivalence, a monotone dose response, a food-service mechanism, a validated latent exposure measure, complete survey-design inference, or replication of proprietary employer hiring and separation results.  It does not infer that a literal vintage merge was used by another study without implementation evidence.

Two supplied-data limitations remain.  Exact BCC public-dashboard occupation membership and tie handling are unavailable, so the CPS and ACS BCC exercises remain labeled approximations.  The authorized full-window extracts omit \texttt{EARNWEEK2}, so the legacy earnings sample ending in March 2023 is not presented as a substitute for a full-window earnings analysis.  Author identity and journal disclosure fields remain marked for completion by the author.

\end{document}
